\section{Experiments}
\label{sec:experiments}

The experiments follow the predictions of Section~\ref{sec:analysis} rather than an average:
we check that the scale family is non-degenerate before training, report downstream quality,
then test Predictions~\ref{pred:levels}--\ref{pred:anchor}, each of which specifies what
should move \emph{and} what should not.

\paragraph{Protocol.}
The student is BERT-base ($110$M, $d_S=768$). We distil from three teachers of different
families and sizes --- BGE-M3, Qwen3-Embedding-0.6B and Qwen3-Embedding-4B --- to check that
the analysis is not an artifact of one teacher's geometry. Teacher embeddings are computed
once over a deduplicated corpus of $13.5$k texts and cached; the teacher is not loaded
during optimisation. We train $5$ epochs with AdamW at $2\times10^{-5}$, batch size $16$,
redrawing candidate sets each epoch. Unless stated otherwise $\mathcal{R}=\{1,2,4\}$,
$k=50$, $\tau_g = 0.05$; all constants are in Appendix~\ref{app:impl}.

\paragraph{Evaluation and baselines.}
Nine benchmarks in three families: linear-probe classification $F_1$ (Bank77, TweetEval,
Emotion), pair classification average precision (MRPC, SciTail, WiC), and STS Spearman
correlation (SICK-R, STS12, STS-B). The split into \emph{level-sensitive} metrics (pair
classification and STS, which score pairs drawn from many different anchors on one common
scale) and \emph{probe-based} metrics (linear probes, which fit their own boundary on the
embeddings and never compare cosines across anchors) is fixed in advance by
Prediction~\ref{pred:levels} and is
the axis along which we read every ablation; three datasets overlap the corpus in domain and
are marked separately. Against the undistilled student and unsupervised SimCSE we compare
five recent embedding distillation methods (CDM, DSKD, Jasper \& Stella, DistillCSE, EMO) and
the strongest prior relational method (TALAS). We report SEED \citep{fang2021seed} explicitly
as the $\Rset=\{0\}$ member of \eqref{eq:family}: it is the reference point against which a
gain is attributable to $\mathcal{R}$ and the temperature bank rather than to relational
distillation as such.

\subsection{The scale family is non-degenerate}
\label{sec:exp:build}

Table~\ref{tab:diag} in Appendix~\ref{app:diag} reports the build-time diagnostics of
Section~\ref{sec:method} for the Qwen3-4B teacher. Sharpness \eqref{eq:diag1} is
bounded away from zero at every scale ($0.65$, $1.72$, $1.53$ nats for $r=1,2,4$), so the
targets carry ranking information rather than a binary neighbour label. The pairwise
divergences are small but non-zero and decrease with $r$ ($0.173$, $0.127$), quantifying the
mixing argument of Section~\ref{sec:setup}: beyond a few lazy steps the scales duplicate one
another, and this diagnostic --- not a downstream sweep --- is what bounds a useful
$\mathcal{R}$. Finally $\JS_\omega = 0.100$ nats: positive, so Corollary~\ref{cor:floor} does
not vacuously certify degeneracy, yet small enough that a tied-readout run would spend most
of its loss on a constant. The truncation residual is numerically zero at $r=1$ and grows to
$8.7\times10^{-2}$ at $r=4$, reported rather than assumed away.

\subsection{Downstream quality}
\label{sec:exp:main}

% Cells marked "--" have not been run yet.
% The \method row for Qwen3-Embedding-4B is a SINGLE SEED; multi-seed pending.
\begin{table}[t]
\centering
\caption{Downstream embedding quality across nine benchmarks. Datasets marked $\star$ are
in-domain with respect to the distillation corpus. Classification is linear-probe macro
$F_1$, pair classification is average precision, STS is Spearman correlation. ``--'' = not
run; the \method row is a single seed.}
\label{tab:main}
\scriptsize
\setlength{\tabcolsep}{2.2pt}
\begin{tabular}{l|ccc|ccc|ccc|cc|c}
\toprule
 & \multicolumn{3}{c|}{\textbf{Classification} ($F_1$)} &
 \multicolumn{3}{c|}{\textbf{Pair Cls.} (AP)} &
 \multicolumn{3}{c|}{\textbf{STS} ($\rho$)} &
 \multicolumn{2}{c|}{\textbf{Domain}} & \\
 & Bank77 & Tweet & Emot.$^{\star}$ &
 MRPC & SciTail & WiC$^{\star}$ &
 SICK & STS12 & STSB$^{\star}$ &
 In & Out & \textbf{Avg} \\
\midrule
Student (no distil.)& 84.86 & 47.79 & \textbf{68.02} & 77.36 & 67.74 & 59.22 & 42.43 & 21.54 & 20.30 & 42.44 & 60.33 & 54.36 \\
SimCSE-unsup       & 89.21 & 69.13 & 53.55 & 82.86 & 76.59 & \textbf{71.64} & 68.19 & 61.68 & 70.36 & 65.18 & 74.61 & 71.47 \\
\midrule
\multicolumn{13}{c}{\textbf{BGE-M3} $\rightarrow$ BERT-base} \\
\midrule
Teacher            & 93.52 & 73.85 & 68.56 & 85.81 & 91.87 & 61.47 & 79.18 & 78.73 & 84.87 & 71.63 & 83.83 & 79.76 \\
CDM                & -- & -- & -- & -- & -- & -- & -- & -- & -- & -- & -- & -- \\
DSKD               & -- & -- & -- & -- & -- & -- & -- & -- & -- & -- & -- & -- \\
Jasper \& Stella   & -- & -- & -- & -- & -- & -- & -- & -- & -- & -- & -- & -- \\
DistillCSE         & -- & -- & -- & -- & -- & -- & -- & -- & -- & -- & -- & -- \\
EMO                & -- & -- & -- & -- & -- & -- & -- & -- & -- & -- & -- & -- \\
TALAS              & -- & -- & -- & -- & -- & -- & -- & -- & -- & -- & -- & -- \\
SEED (= $r{=}0$ only)& -- & -- & -- & -- & -- & -- & -- & -- & -- & -- & -- & -- \\
\textbf{\method}   & -- & -- & -- & -- & -- & -- & -- & -- & -- & -- & -- & -- \\
\midrule
\multicolumn{13}{c}{\textbf{Qwen3-Embedding-0.6B} $\rightarrow$ BERT-base} \\
\midrule
Teacher            & 93.21 & 71.33 & 66.53 & 83.37 & 90.00 & 66.85 & 80.64 & 77.11 & 84.60 & 72.66 & 82.61 & 79.29 \\
CDM                & -- & -- & -- & -- & -- & -- & -- & -- & -- & -- & -- & -- \\
DSKD               & -- & -- & -- & -- & -- & -- & -- & -- & -- & -- & -- & -- \\
Jasper \& Stella   & -- & -- & -- & -- & -- & -- & -- & -- & -- & -- & -- & -- \\
DistillCSE         & -- & -- & -- & -- & -- & -- & -- & -- & -- & -- & -- & -- \\
EMO                & -- & -- & -- & -- & -- & -- & -- & -- & -- & -- & -- & -- \\
TALAS              & -- & -- & -- & -- & -- & -- & -- & -- & -- & -- & -- & -- \\
SEED (= $r{=}0$ only)& -- & -- & -- & -- & -- & -- & -- & -- & -- & -- & -- & -- \\
\textbf{\method}   & -- & -- & -- & -- & -- & -- & -- & -- & -- & -- & -- & -- \\
\midrule
\multicolumn{13}{c}{\textbf{Qwen3-Embedding-4B} $\rightarrow$ BERT-base} \\
\midrule
Teacher            & 93.63 & 72.22 & 68.94 & 84.79 & 91.77 & 66.80 & 83.43 & 82.55 & 86.87 & 74.20 & 84.73 & 81.22 \\
CDM                & 89.75 & 70.58 & 53.11 & 84.83 & 75.19 & 71.06 & 69.37 & 68.48 & 73.92 & 66.03 & 76.37 & 72.92 \\
DSKD               & 89.66 & 69.93 & 54.51 & 83.13 & 77.95 & \underline{71.53} & 68.65 & 63.37 & 71.42 & 65.82 & 75.45 & 72.24 \\
Jasper \& Stella   & 89.71 & \textbf{73.83} & 65.98 & \underline{85.33} & 80.99 & 70.50 & 73.96 & 67.90 & 75.65 & 70.71 & 78.62 & 75.98 \\
DistillCSE         & 90.94 & 70.90 & 60.80 & 82.86 & 78.98 & 68.63 & 75.12 & \underline{72.61} & 77.08 & 68.84 & 78.57 & 75.32 \\
EMO                & 90.78 & 73.54 & 67.48 & 84.46 & 81.57 & 69.76 & 75.66 & 68.83 & 77.17 & \underline{71.47} & 79.14 & 76.58 \\
TALAS              & \textbf{92.34} & \underline{73.74} & 64.41 & \textbf{85.97} & \underline{82.66} & 68.65 & \textbf{78.38} & 72.27 & \underline{79.98} & 71.01 & \underline{80.89} & \underline{77.60} \\
SEED (= $r{=}0$ only)& -- & -- & -- & -- & -- & -- & -- & -- & -- & -- & -- & -- \\
\textbf{\method}   & \underline{91.43} & 72.84 & \underline{67.74} & 84.15 & \textbf{85.38} & 68.65 & \underline{77.81} & \textbf{74.27} & \textbf{82.99} & \textbf{73.13} & \textbf{80.98} & \textbf{78.36} \\
\bottomrule
\end{tabular}
\end{table}


For the Qwen3-Embedding-4B teacher, \method reaches $78.36$ average against $77.60$ for
TALAS, the strongest prior method. The gain is not uniform, but neither is it as cleanly
differential as Prediction~\ref{pred:levels} asks. The three largest moves are
level-sensitive and positive --- STS-B ($+3.01$), STS12 ($+2.00$), SciTail ($+2.72$) --- yet
MRPC ($-1.82$) and SICK-R ($-0.57$) move the other way, so the level-sensitive average rises
by only $+0.89$. The probe-based average rises by $+0.51$, and that figure is carried
entirely by Emotion ($+3.33$) against declines on Bank77 ($-0.91$) and TweetEval ($-0.90$).
A differential of $+0.89$ against $+0.51$ points in the predicted direction and is, at one
seed, too small to be evidence for it. We therefore do not read
Prediction~\ref{pred:levels} off this table at all: a comparison against a different method
varies the objective, the targets and the optimisation together, whereas the ablations of
Section~\ref{sec:exp:abl} vary one component against a fixed baseline and are where the
prediction is actually tested. In-domain and out-of-domain averages both improve over TALAS
($+2.12$ and $+0.09$), so the effect is not a corpus-overlap artifact, though the
out-of-domain margin is small.

The remaining teacher blocks and the SEED row are not yet run and are shown as ``--'', which
leaves three things unestablished: the \method row is a \emph{single seed}; the comparison
to SEED --- the ablation separating our contribution from the relational baseline it builds
on --- is pending; and the claim that the analysis is teacher-independent rests on one teacher.

\subsection{Testing the analysis}
\label{sec:exp:abl}

Each result of Section~\ref{sec:analysis} names a quantity that a single readout leaves
free, so each can be tested twice: \emph{directly}, by measuring that quantity inside the
model, and \emph{indirectly}, by the downstream damage its absence causes.
Table~\ref{tab:mech} reports the first and Table~\ref{tab:ablation} the second. We weight
the direct measurements more heavily. A downstream average is a noisy, many-step
consequence of the mechanism a proposition describes, whereas the excess over the floor,
the level gap, the off-graph repulsion rate and the anchor gradient share are the
propositions themselves written as numbers --- and a mechanism that fails to appear there
is not rescued by a favourable benchmark average.

% Mechanism measurements. Nothing run yet.
\begin{table}[t]
\centering
\caption{Mechanism measurements: for each result of Section~\ref{sec:analysis}, the quantity
that result is about, measured inside the model rather than through a downstream proxy.
\emph{Excess} is $\mathcal{L}-\JS_\omega$ at convergence, which Corollary~\ref{cor:floor}
forces to $\geq 0$ under a single readout and only bounds below under a bank. \emph{Level
gap} is the mean within-anchor Spearman correlation between student and teacher cosines
minus the same correlation computed on all pairs pooled across anchors; it is exactly the
discrepancy the per-anchor shift of Proposition~\ref{prop:levels}(a) creates.
\emph{Off-graph repulsion} is the fraction of teacher-near, off-pool pairs whose student
cosine decreases over training, which \eqref{eq:zerograd} drives to $1$ absent a
full-support target. \emph{Anchor share} is
$\Vert\nabla_\theta\mathcal{L}_{\mathrm{anc}}\Vert/\Vert\nabla_\theta\mathcal{L}\Vert$ at the
last epoch. ``--'' = not run.}
\label{tab:mech}
\small
\setlength{\tabcolsep}{5pt}
\begin{tabular}{lcccc}
\toprule
 & Excess & Level gap & Off-graph & Anchor \\
Variant & $\mathcal{L}-\JS_\omega$ & (Prop.~\ref{prop:levels}) & repulsion & share \\
 & (Cor.~\ref{cor:floor}) &  & (Prop.~\ref{prop:support}) & (Lem.~\ref{lem:gauge}) \\
\midrule
\multicolumn{5}{l}{\emph{Predicted}}\\
Single readout & $\to 0^{+}$ & large & --- & --- \\
\method (full) & $<0$ & small & $\approx$ base rate & --- \\
No full-support target & --- & --- & $\to 1$ & --- \\
With $\mathcal{L}_{\mathrm{anc}}$ & --- & unchanged & --- & $\to 0$ \\
\midrule
\multicolumn{5}{l}{\emph{Measured}}\\
\textbf{\method (full)} & -- & -- & -- & n/a \\
Tied readout, $\tau_r\equiv\tau$ & -- & -- & -- & n/a \\
SEED, $\Rset=\{0\}$ & n/a & -- & -- & n/a \\
Drop $r=0$ (graph only) & -- & -- & -- & n/a \\
Full $+\,\mathcal{L}_{\mathrm{anc}}$, $W_a$ free & -- & -- & -- & -- \\
\bottomrule
\end{tabular}
\end{table}


\paragraph{The statistic the predictions actually make.}
Predictions~\ref{pred:levels} and~\ref{pred:support} do not say ``the variant is worse'';
they say it is worse \emph{on level-sensitive metrics specifically}. The corresponding
statistic is the excess damage
$D = \Delta_{\mathrm{Probe}} - \Delta_{\mathrm{Level}}$ relative to the full method, and a
variant that degrades both families equally has $D\approx 0$ and refutes the prediction as
surely as one that degrades neither. $D$ is a difference of two noisy averages, so it is
uninterpretable without a scale: we run the full method at three seeds and report its spread
in the same units, and count no $D$ that does not exceed it. This is also why we do not read
these predictions off Table~\ref{tab:main}: comparing against a different method varies the
objective, the targets and the optimisation at once, whereas every row below changes one
component against a fixed baseline.

% Downstream ablations. Only the reference rows are filled (single seed); all variants "--".
\begin{table}[t]
\centering
\caption{Downstream ablations, grouped by the result each isolates. \emph{Level-sens.}
averages the six metrics that score pairs drawn from many different anchors on one common
scale (MRPC, SciTail, WiC average precision; SICK-R, STS12, STS-B Spearman $\rho$), which
depend on the per-anchor level Proposition~\ref{prop:levels}(a) leaves free;
\emph{Probe} averages the three linear probes (Bank77, TweetEval, Emotion), which fit their
own boundary on the embeddings and never compare cosines across anchors. Since
Predictions~\ref{pred:levels}--\ref{pred:anchor} are differential, the test statistic is
$D = \Delta_{\mathrm{Probe}} - \Delta_{\mathrm{Level}}$, the excess damage a variant does to
the level-sensitive family; $D>0$ supports a prediction, $D\approx0$ refutes it whether or
not the variant is worse overall. \emph{Pred.} fixes the expected sign before the run
($++$ large, $+$ positive, $0$ none). $D$ counts only if it exceeds the seed spread in the
reference block. ``--'' = not run.}
\label{tab:ablation}
\small
\setlength{\tabcolsep}{4pt}
\begin{tabular}{llcccc}
\toprule
Variant & Isolates & Pred. & Level-sens. & Probe & $D$ \\
\midrule
\multicolumn{6}{l}{\emph{Reference points}}\\
\textbf{\method (full)} & --- & --- & \textbf{78.88} & \textbf{77.34} & --- \\
\quad seed spread, $n=3$ & noise floor for $D$ & --- & -- & -- & -- \\
SEED, $\Rset=\{0\}$ & no bank, no diffusion & $+$ & -- & -- & -- \\
TALAS (strongest prior) & different objective & n/a & 77.98 & 76.83 & n/a \\
\midrule
\multicolumn{6}{l}{\emph{(A) Prop.~\ref{prop:collapse}: one readout collapses the family onto its mixture}}\\
Tied readout, $\tau_r\equiv\tau$ & Pred.~\ref{pred:levels} & $++$ & -- & -- & -- \\
\quad vs.\ single scale on $\bar p$ & Eq.~\eqref{eq:collapse} is exact & $0$ & -- & -- & -- \\
Mixture-proportional sampler & Cond.~\ref{cond:cover}: collapse by sampling & $+$ & -- & -- & -- \\
\midrule
\multicolumn{6}{l}{\emph{(B) Prop.~\ref{prop:levels}: one readout leaves the per-anchor level free}}\\
Two distinct $\tau$, $\Rset$ fixed & part (b), minimal case & $+$ & -- & -- & -- \\
$\tau_r$ decreasing in $r$ & distinctness vs.\ assignment & $0$ & -- & -- & -- \\
Per-anchor columns, no $\mathcal{C}_B$ & levels across anchors & $+$ & -- & -- & -- \\
\midrule
\multicolumn{6}{l}{\emph{(C) Prop.~\ref{prop:support}: truncated targets repel off the graph}}\\
Drop $r=0$ (graph only) & Pred.~\ref{pred:support} & $++$ & -- & -- & -- \\
\quad $r{=}0$ restricted to $\mathcal{P}_i$ & support alone, $\omega_0$ and mass held & $++$ & -- & -- & -- \\
\quad $r{=}0$ replaced by uniform & support without the teacher's mass & $+$ & -- & -- & -- \\
\midrule
\multicolumn{6}{l}{\emph{(D) Lem.~\ref{lem:gauge}: a pointwise anchor cannot supply the missing quantity}}\\
Tied readout $+\,\mathcal{L}_{\mathrm{anc}}$ & can it substitute for the bank? & $++$ & -- & -- & -- \\
Full $+\,\mathcal{L}_{\mathrm{anc}}$, $W_a$ free & Pred.~\ref{pred:anchor}, a null & $0$ & -- & -- & -- \\
Full $+\,\mathcal{L}_{\mathrm{anc}}$, $W_a$ orth. & control: $GL(d_S)$ removed & $0$ & -- & -- & -- \\
\bottomrule
\end{tabular}
\end{table}


\paragraph{(A) The collapse.}
The tied-readout variant holds the target family and the weights fixed and changes only
whether the student side varies with $r$. Here the theory admits a check sharper than any
metric comparison: by \eqref{eq:collapse} the tied run and a single-scale run on $\bar p$
are the \emph{same optimisation problem}, so we compare their per-step losses and gradients
directly and report the maximum discrepancy, which should be at numerical precision. A
downstream disagreement between those two rows indicts the implementation, not the analysis.
The mixture-proportional sampler then shows the same failure entering through the candidate
draw rather than the objective: it violates Condition~\ref{cond:cover} by starving the
sharpest scale of scored columns, and should reproduce part of the tied-readout damage even
though \eqref{eq:objective} is untouched.

\paragraph{(B) The levels.}
The level gap of Table~\ref{tab:mech} is the direct test: if within-anchor agreement with
the teacher is high while pooled agreement is much lower, the student has learned the
ranking and not the levels, which is Proposition~\ref{prop:levels}(a) verbatim. Two rows
then probe how much of the bank is needed. Using two distinct temperatures while keeping
$\Rset$ unchanged tests the minimal case of part (b) without also removing scales --- a
variant that shrank $\Rset$ would confound the size of the bank with the size of the family
--- and making $\tau_r$ decrease with $r$ tests whether only the \emph{distinctness} of the
temperatures is load-bearing, as the analysis claims, or whether the monotone assignment we
adopt is doing work the analysis does not account for. Dropping the shared column set
restores the per-anchor freedom that part (b) removes only when anchors compete over common
columns.

\paragraph{(C) The support.}
Dropping $r=0$ leaves no full-support target, so by \eqref{eq:zerograd} every off-pool
candidate is pushed away at a rate independent of its teacher similarity; the off-graph
repulsion rate should approach $1$, and its downstream shadow should fall in the
level-sensitive column. That variant alone is confounded, however: removing $r=0$ also
renormalises $\omega$ over the remaining scales and discards the only dense teacher signal.
We therefore add a control that keeps $\omega_0$ and the teacher's mass and changes support
alone, restricting $p_{i,0}$ to the pool $\mathcal{P}_i$. If the damage is a support effect
this row reproduces most of it; if it does not, what mattered was the dense signal rather
than Proposition~\ref{prop:support}. Replacing $r=0$ by uniform smoothing completes the
dissection, separating \emph{having} full support from \emph{having the teacher's} mass on
it, and should recover only part of the gap.

\paragraph{(D) The pointwise anchor.}
The question is substitution, not addition, so the load-bearing row adds
$\mathcal{L}_{\mathrm{anc}}$ to a student that still lacks the bank: if a pointwise term
supplied calibration, tied readout $+\,\mathcal{L}_{\mathrm{anc}}$ would recover the
level-sensitive column, and Lemma~\ref{lem:gauge} says it cannot. Adding it to the full
method is the complementary null --- worth committing to in advance --- and the anchor
gradient share measures the mechanism: the term should fit the best linear predictor of
$t_i$ from $s_i$ early and then contribute almost nothing. The orthogonal-$W_a$ row is the
control that distinguishes the two available explanations: if constraining $W_a$ helps where
a free $W_a$ does not, the failure was the $GL(d_S)$ gauge rather than pointwise supervision
as such.
